\section{Complete finite classification}\label{sec:classification}
We now retain only the necessary information proved in the preceding
section. The resulting classes are deliberately larger than the set of
winning strategies. Excluding all of them will therefore exclude every
winning strategy, even though some listed vectors may be unrealizable
for reasons not used in the classification.
\subsection{Count representatives}
For an integer $n$, let $\calA_n$ be the set of integer vectors
$a=(a_0,\ldots,a_4)$ satisfying
\begin{align}
0\leq a_j\leq16,\qquad&\sum_{j=0}^4a_j=n,\label{eq:A1}\\
a_j+a_{j+1}&\geq6,\label{eq:A2}\\
a_j+a_{j+1}=6\quad&\Longrightarrow\quad
(a_j,a_{j+1})\in\{(0,6),(2,4),(3,3),(4,2),(6,0)\},\label{eq:A3}\\
a_{j-1}+a_j=a_j+a_{j+1}=6\quad&\Longrightarrow\quad
a_j\in\{0,4\}.\label{eq:A4}
\end{align}
All indices in these conditions are modulo~5. These conditions are
necessary, not asserted to be sufficient for a strategy.

Let $D_5$ act on these vectors by cyclic shifts and reversal. Choose the
representative that minimizes the ordered pair
\begin{equation}\label{eq:canonical}
\bigl(a_0+a_1,(a_0,a_1,a_2,a_3,a_4)\bigr)
\end{equation}
lexicographically, and denote the set of representatives by $\calR_n$.
The first component of~\eqref{eq:canonical} is an ordering convention
that places a least pair sum first; it is not an additional restriction.
After one representative has been selected from each orbit, the saved
list is sorted lexicographically by the five-tuple itself. This
serialization order is distinct from the key used to select a
representative; set equality and absence of duplicates determine coverage.

\begin{lemma}[Validity of the graph reduction]\label{lem:dihedral}
Every winning strategy with $S(z)=n$ has an equivalent strategy whose
count vector in the order $(1,3,5,2,4)$ belongs to $\calR_n$.
\end{lemma}
\begin{proof}
The count vector belongs to $\calA_n$ by the results of
Section~\ref{sec:local}. A rotation of the vertices transports the five
functions in their cyclic order. A reflection transports the functions
and interchanges their left and right arguments. In either case proper
colorings are carried bijectively to proper colorings, and correct guesses
are preserved. On the matrices, reflection includes transposition, which
preserves their weights. The action on the order $(1,3,5,2,4)$ is again all
cyclic shifts and reversals: multiplication by 2 is invertible modulo~5.
Thus an automorphism sends the count vector to its representative without
requiring any invariance of the strategy itself.
\end{proof}

For clarity,
use vertex labels in $\mathbb Z/5\mathbb Z$ temporarily. A rotation
$h(i)=i+d$ transports rules by $g_{h(i)}(u,v)=f_i(u,v)$; a reflection
$h(i)=d-i$ transports them by $g_{h(i)}(u,v)=f_i(v,u)$. For
$c'_{h(i)}=c_i$, both formulas give
$g_{h(i)}(c'_{h(i)-1},c'_{h(i)+1})=f_i(c_{i-1},c_{i+1})$.
Thus each individual correct guess is preserved. In the count order
$i=1+2j$, the induced actions are $j\mapsto j+k$ and $j\mapsto k-j$.

There are two elementary ways to produce the entire list. One enumerates
$(a_0,a_1,a_2,a_3)\in\{0,\ldots,16\}^4$, sets
$a_4=n-a_0-a_1-a_2-a_3$, and checks~\eqref{eq:A1}--\eqref{eq:A4}. A
different enumeration uses slack variables
\begin{equation}\label{eq:slack}
e_j=a_j+a_{j+1}-6,\qquad e_j\geq0,\qquad\sum_je_j=2n-30.
\end{equation}
Write $b_j=6+e_j$. The equations $a_j+a_{j+1}=b_j$ have the unique solution
\begin{equation}\label{eq:inverse}
2a_0=b_0-b_1+b_2-b_3+b_4,\qquad
a_{j+1}=b_j-a_j\quad(0\leq j<4).
\end{equation}
Enumerating all compositions of $2n-30$ into five nonnegative parts,
applying~\eqref{eq:inverse}, and testing~\eqref{eq:A1}--\eqref{eq:A4}
therefore gives exactly the same set $\calA_n$. An odd numerator in
\eqref{eq:inverse} is rejected. The last cyclic equation follows from the
formula for $a_0$; the total sum is also checked explicitly in the audit.
This establishes completeness of both enumeration schemes.

For example, the list at weight 17 consists of
\[
\calR_{17}=\{(2,4,3,3,5),\ (2,4,4,2,5),\ (3,3,4,3,4)\}.
\]
The middle vector specifies player counts $(2,2,4,5,4)$ in the ordinary
order $(1,2,3,4,5)$, but fixes none of the locations of the counted entries.
It is also the count vector of the local witness in Appendix~\ref{app:witness}.

\subsection{Color relabeling and 51 diagonal patterns}
A permutation $\sigma$ of $Q$ acts on strategies by
\begin{equation}\label{eq:color-action}
f_i^\sigma(a,b)=
\sigma\bigl(f_i(\sigma^{-1}(a),\sigma^{-1}(b))\bigr).
\end{equation}
It preserves properness and correctness of each guess. After making the
chosen color equal to 5, permutations of the remaining four colors preserve
every $s_i(5)$. In particular, since $f_1(5,5)\ne5$, we may impose
\begin{equation}\label{eq:one-unit}
f_1(5,5)=1.
\end{equation}
No second normalization involving a different observation is imposed in
the count-branch inputs.

\begin{lemma}[Diagonal-pattern reduction]\label{lem:patterns}
Let $t_i=f_i(5,5)\in\{1,2,3,4\}$. Up to a common permutation of these
four colors, $(t_1,\ldots,t_5)$ is one of the 51 words satisfying
\begin{equation}\label{eq:rgs}
t_1=1,\qquad
1\leq t_j\leq\min\{4,1+\max_{\ell<j}t_\ell\}\quad(2\leq j\leq5).
\end{equation}
This normalization preserves the count representative already chosen.
\end{lemma}
\begin{proof}
Assign the label 1 to the first color encountered in the word, label 2 to
the next previously unseen color, and so forth. Extend this assignment
arbitrarily to a permutation of all four colors. The resulting word
satisfies~\eqref{eq:rgs}. Conversely, its equality classes determine it
uniquely. These words correspond to partitions of five labeled positions
into at most four nonempty blocks. The recurrence
$\Stirling{n}{k}=k\Stirling{n-1}{k}+\Stirling{n-1}{k-1}$ gives
\[
\Stirling{5}{1}+\Stirling{5}{2}+\Stirling{5}{3}+\Stirling{5}{4}
=1+15+25+10=51.
\]
The permutation fixes color 5, so it only permutes the 16 observations
inside each $M_i^5$ and preserves its weight. If~\eqref{eq:one-unit}
already holds, the first-occurrence relabeling fixes its answer~1.
\end{proof}

\begin{proposition}[Coverage of all five-color candidates]\label{prop:coverage}
If a five-color winning strategy exists, there are
$n\in\{17,\ldots,25\}$ and $a\in\calR_n$ such that an equivalent
normalized winning strategy has count vector $a$ and satisfies
\eqref{eq:one-unit}. For any such count branch selected for refinement,
at least one of the 51 constraints~\eqref{eq:rgs} still contains an
equivalent winning strategy.
\end{proposition}
\begin{proof}
Normalize by Lemma~\ref{lem:normalization}, choose a color with total at
most 25 by Lemma~\ref{lem:total}, and relabel it as~5.
Proposition~\ref{prop:seventeen} supplies the lower bound 17. Apply
Lemma~\ref{lem:dihedral}, then a color permutation fixing 5 to obtain
\eqref{eq:one-unit}. If the branch is refined, apply
Lemma~\ref{lem:patterns}. Each step preserves winning and all count
conditions needed by the next step.
\end{proof}

The 51 subcases need not exhaust all labeled diagonal words as literal
assignments. They cover them up to the explicitly proved color action.
This distinction is essential: the proof does not impose a symmetry
condition on the unknown functions.

The two reductions are applied in a fixed order. First a vertex
 automorphism selects the count representative; next a permutation of
 the four non-special colors selects a diagonal representative. The
 second operation leaves every count unchanged, so it cannot undo the
 first. There is no requirement that the diagonal tuple also be minimal
 under vertex automorphisms.

